% Preamble
\documentclass[12pt,aspectratio=169]{beamer}
\usepackage{graphicx}
\usepackage{xcolor}
\usepackage{fontspec}
\usepackage{tikz}

\setsansfont{HTWBerlin Office}

\usetheme{default}

\definecolor{htwgreen}{RGB}{118,185,0}
\definecolor{darkgreen}{RGB}{30,130,0}
\definecolor{darkred}{RGB}{175,0,0}
\setbeamercolor*{normal text}{fg=white,bg=htwgreen}
\setbeamercolor*{structure}{use=normal text}

\setbeamertemplate{frametitle}{
    \vspace{0.7cm}
    \textbf{\insertframetitle}

    \large\insertframesubtitle
}
\setbeamertemplate{navigation symbols}{}

\logo{\includegraphics[width=1.5cm]{resources/fsr4}}

% Document
\begin{document}

    \begin{frame}
        \huge
        \textbf{Schön, dass ihr da seid :)}

        \medskip

        \normalsize
        \textbf{Herzlich willkommen!}

        \bigskip

        \textbf{Hier ein paar Regeln:}

        \begin{itemize}
            \item Wortmeldungen bitte vor allem in den Chat
            \item Um Störgeräusche zu vermeiden, schaltet eure Mikrofone
            bitte nur ein, wenn ihr das Wort erteilt bekommen habt.
        \end{itemize}
    \end{frame}

    \begin{frame}[squeeze]
        \frametitle{Tagesordnung}

        \begin{columns}

            \begin{column}{0.5\textwidth}
                \begin{enumerate}
                    <BLOCK>for top in tops_left</BLOCK>
                        \item <VAR>top.name</VAR>
                        <BLOCK>if top.sub</BLOCK>
                            \begin{enumerate}
                                <BLOCK>for top in top.sub</BLOCK>
                                    \item <VAR>top.name</VAR>
                                <BLOCK>endfor</BLOCK>
                            \end{enumerate}
                        <BLOCK>endif</BLOCK>
                    <BLOCK>endfor</BLOCK>
                \end{enumerate}
            \end{column}

            \vrule width 1pt

            \begin{column}{0.5\textwidth}
                \begin{enumerate}
                    \setcounter{enumi}{<VAR>tops_left | count</VAR>}
                    <BLOCK>for top in tops_right</BLOCK>
                        \item <VAR>top.name</VAR>
                        <BLOCK>if top.sub</BLOCK>
                            \begin{enumerate}
                                <BLOCK>for top in top.sub</BLOCK>
                                    \item <VAR>top.name</VAR>
                                <BLOCK>endfor</BLOCK>
                            \end{enumerate}
                        <BLOCK>endif</BLOCK>
                    <BLOCK>endfor</BLOCK>
                \end{enumerate}
            \end{column}

        \end{columns}
    \end{frame}

    \begin{frame}
        \frametitle{Abstimmung}
        \framesubtitle{Bitte die Maus in den zutreffenden Kasten bewegen}

        (nur für Mitglieder des Fachschaftsrats 4)

        \bigskip

        \begin{columns}

            \begin{column}{0.3\textwidth}
                \centering
                JA
                \bigskip

                \begin{tikzpicture}
                    \draw[darkgreen,fill=darkgreen] (0,0) rectangle (4,4);
                \end{tikzpicture}
            \end{column}

            \begin{column}{0.3\textwidth}
                \centering
                NEIN
                \bigskip

                \begin{tikzpicture}
                    \draw[darkred,fill=darkred] (0,0) rectangle (4,4);
                \end{tikzpicture}
            \end{column}

            \begin{column}{0.3\textwidth}
                \centering
                ENTHALTUNG
                \bigskip

                \begin{tikzpicture}
                    \draw[lightgray,fill=lightgray] (0,0) rectangle (4,4);
                \end{tikzpicture}
            \end{column}

        \end{columns}
    \end{frame}

    \begin{frame}
        \frametitle{Meinungsbild}
        \framesubtitle{Bitte die Maus in den zutreffenden Kasten bewegen}

        (Hier dürfen/müssen alle mitmachen :))

        \bigskip

        \begin{columns}

            \begin{column}{0.5\textwidth}
                \centering
                JA
                \bigskip

                \begin{tikzpicture}
                    \draw[darkgreen,fill=darkgreen] (0,0) rectangle (4,4);
                \end{tikzpicture}
            \end{column}

            \begin{column}{0.5\textwidth}
                \centering
                NEIN
                \bigskip

                \begin{tikzpicture}
                    \draw[darkred,fill=darkred] (0,0) rectangle (4,4);
                \end{tikzpicture}
            \end{column}

        \end{columns}
    \end{frame}

\end{document}
